% page 477 du fac-simile (image p-476 du scan)
% bloc 162.6 x 233.3 mm ; marge gauche 39.4 mm
\pgc{17.780mm}{0.000mm}{
\l{\cel{49}469}
\l{}
\l{\sou{prozo}. Formo di la pens-expresuro qua ne submisesas a la mezuro}
\l{\cel{3}ed a la ritmo di la verso. -\cel{2}DEFIRS.}
\l{}
\l{\sou{prozelito}. Nova adepto di religio. - DEFIRS.}
\l{}
\l{\sou{prozodio}. I. (\sou{gram.}) Pronunco regulkonforma di vorto, konforme a}
\l{\cel{3}la acentizo. - II. Reguli pri la quanto de silabi e la mezuro}
\l{\cel{3}di la versi diversa. - DEFIRS.}
\l{}
\l{\sou{prozopopeo.}\cel{2}(\sou{retor.}) Figuro qua konsistas ye egardigar kom}
\l{\cel{3}vivanta, kozin, personin mortinta od absenta. - DEFIS.}
\l{}
\l{\sou{pruo.}\cel{2}Avana parto di navo, di batelo. - EFIS.}
\l{}
\l{\sou{pruda}. Qua afektacas vertuo rigoroza. - DEF.}
\l{}
\l{\sou{prudenta}. Qua praktikas ta sajeso qua regulas la konduto ed igas}
\l{\cel{3}posibla la evito di kulpi, di erori. - EFIS.}
\l{}
\l{\sou{pruinar}. (\sou{netrans}.)Kovreskar per roso konjelita, o per nebulo}
\l{\cel{3}konjelita, lor la nokti kolda di autuno o di printempo.}
\l{}
\l{}
\l{\sou{pruno}. (\sou{bot.)}\cel{2}Frukto kernoza, kun shelo glata. - L.\cel{1}\sou{prunus}.EFI.}
\l{}
\l{\sou{prunelo}. (\sou{bot.}) Pruno sovaja, nigra, di qua la saporo esas akra.}
\l{\cel{3}- L.\cel{1}\sou{prunus spinosa}. - FI.}
\l{}
\l{\sou{pruntar}. (\sou{trans.,}\cel{1}de).I. Obtenar ke ulu prestas. - II. Tirar}
\l{\cel{3}(ulo) de kunhomo. - III. Rekursar ad helpo de altru. - F.}
\l{}
\l{\sou{pruritar}. (\sou{netrans.}) Tale priklar ke divenas necesega skrapar la}
\l{\cel{3}pelo. - EFIS.}
\l{}
\l{\sou{pruvar}. (\sou{trans.)}. Vidigar kom exakta, per rezono o per fakti.}
\l{\cel{3}- DEFIS.}
\l{}
\l{\sou{psalmo}. (\sou{religio krist}.) Kantiko religiala. - DEIFS.}
\l{}
\l{\sou{psalmodiar}. (\sou{trans.}) I. (\sou{kirko).}\cel{1}Kantar psalmi. -- II. (\sou{metaf.})}
\l{\cel{3}Recitar monotone. - DEFIS.}
\l{}
\l{\sou{psalterio}. (\sou{epoki antiqua}). Muzik-instrumento kordoza, quan onu}
\l{\cel{3}pleas per plektro. - DEFIRS.}
\l{}
\l{\sou{pseudo}- .Prefixo qua havas kom senco \textquotedbl{}ne-exakta, ne-reala\textquotedbl{}.}
\l{}
\l{\sou{pseudonimo}. Qua publikigas (o publikigesas) kun autor-nomo ne-}
\l{\cel{3}exakta.\cel{2}-\cel{1}\sou{Pseudonimo}\cel{1}:\cel{2}Nomo ne-exakta. - DEFIRS.}
\l{}
\l{\sou{pseudosfero}. (\sou{geom.}) Surfaco di rotaco, produktita da la trakto-}
\l{\cel{3}rio qua rotacas cirkum sua asimptoto. -}
}
